
\documentclass{article}
\usepackage{colortbl}
\usepackage{makecell}
\usepackage{multirow}
\usepackage{supertabular}

\begin{document}

\newcounter{utterance}

\centering \large Interaction Transcript for game `hot\_air\_balloon', experiment `air\_balloon\_survival\_en\_reasoning off\_hard', episode 4 with gpt{-}5{-}2025{-}08{-}07{-}no{-}reasoning{-}t1.0.
\vspace{24pt}

{ \footnotesize  \setcounter{utterance}{1}
\setlength{\tabcolsep}{0pt}
\begin{supertabular}{c@{$\;$}|p{.15\linewidth}@{}p{.15\linewidth}p{.15\linewidth}p{.15\linewidth}p{.15\linewidth}p{.15\linewidth}}
    \# & $\;$A & \multicolumn{4}{c}{Game Master} & $\;\:$B\\
    \hline

    \theutterance \stepcounter{utterance}  
    & & \multicolumn{4}{p{0.6\linewidth}}{
        \cellcolor[rgb]{0.9,0.9,0.9}{
            \makecell[{{p{\linewidth}}}]{
                \texttt{\tiny{[P1$\langle$GM]}}
                \texttt{Stai partecipando a un gioco di negoziazione collaborativa.} \\
\\ 
\texttt{Insieme a un altro partecipante, dovete accordarvi su un unico set di oggetti da conservare. Ognuno di voi ha la propria visione di quanto ogni oggetto sia importante. Non sai come l'altro partecipante valuta gli oggetti. Inoltre, ti viene fornito lo sforzo che ogni oggetto richiede.} \\
\texttt{Potete accordarvi su un set solo se lo sforzo totale degli oggetti selezionati non supera un limite condiviso:} \\
\\ 
\texttt{LIMITE = 2984} \\
\\ 
\texttt{Ecco i valori di sforzo dei singoli oggetti:} \\
\\ 
\texttt{Sforzo oggetto = \{"A21": 350, "B93": 276, "A11": 616, "B33": 737, "C08": 531, "B39": 389, "B96": 24, "C85": 125, "A82": 338, "B31": 356, "B20": 143, "A60": 117, "B09": 257, "C81": 921, "B21": 789\}} \\
\\ 
\texttt{Ecco la tua visione personale sull'importanza di ogni oggetto:} \\
\\ 
\texttt{Valori di importanza dell'oggetto = \{"A21": 138, "B93": 583, "A11": 868, "B33": 822, "C08": 783, "B39": 65, "B96": 262, "C85": 121, "A82": 508, "B31": 780, "B20": 461, "A60": 484, "B09": 668, "C81": 389, "B21": 808\}} \\
\\ 
\texttt{Obiettivo:} \\
\\ 
\texttt{Il tuo obiettivo è negoziare un set condiviso di oggetti che ti avvantaggi il più possibile (cioè, massimizzi l'importanza totale per TE), rimanendo entro il LIMITEE. Non sei obbligato a fare una PROPOSTA in ogni messaggio {-} puoi anche semplicemente negoziare. Tutte le tattiche sono permesse!} \\
\\ 
\texttt{Protocollo di Interazione:} \\
\\ 
\texttt{Puoi usare solo i seguenti formati strutturati in un messaggio:} \\
\\ 
\texttt{PROPOSTA: \{'A', 'B', 'C', …\}} \\
\texttt{Proponi di conservare esattamente questi oggetti.} \\
\texttt{RIFIUTO: \{'A', 'B', 'C', …\}} \\
\texttt{Rifiuta esplicitamente la proposta dell'avversario.} \\
\texttt{ARGOMENTAZIONE: \{'...'\}} \\
\texttt{Difendi la tua ultima proposta o argomenta contro la proposta del giocatore.} \\
\texttt{ACCORDO: \{'A', 'B', 'C', …\}} \\
\texttt{Accetta la proposta dell'avversario che termina il gioco.} \\
\\ 
\\ 
\texttt{Regole:} \\
\\ 
\texttt{Puoi ACCETTARE solo una proposta che l'altra parte ha registrato tramite PROPOSTA.} \\
\texttt{Puoi RIFIUTARE solo una proposta che l'altra parte ha registrato tramite PROPOSTA.} \\
\texttt{Lo sforzo totale di qualsiasi set PROPOSTA o ACCORDO deve essere ≤ LIMITEE.} \\
\texttt{NON rivelare i tuoi punteggi di importanza nascosti.} \\
\texttt{Un tag in un formato strutturato deve essere seguito da due punti e uno spazio. L'argomento deve essere un set python contenente 0 o più stringhe.} \\
\texttt{Quindi, deve essere della forma TAG: \{...\}} \\
\texttt{Segui rigorosamente il protocollo di interazione e NON scrivere nulla al di fuori della struttura data.} \\
\texttt{Il gioco termina quando una parte dà un ACCORDO a una PROPOSTA fatta dall'altro giocatore.} \\
\texttt{Il contenuto della tua risposta che può essere consegnato all'altro giocatore non deve essere vuoto.} \\
\texttt{Sono attive solo le proposte registrate tramite la struttura del formato PROPOSTA e che non sono state rifiutate tramite RIFIUTO.} \\
\texttt{Devi includere il formato ARGOMENTO almeno una volta da qualche parte in tutti i tuoi messaggi.} \\
            }
        }
    }
    & \\ \\

    \theutterance \stepcounter{utterance}  
    & \multicolumn{4}{p{0.6\linewidth}}{
        \cellcolor[rgb]{0.9,0.9,0.9}{
            \makecell[{{p{\linewidth}}}]{
                \texttt{\tiny{[P1$\rangle$GM]}}
                \texttt{PROPOSTA: \{'A11', 'B33', 'B31', 'B93', 'B09', 'A60'\}} \\
\\ 
\texttt{ARGOMENTAZIONE: \{'Propongo questo set perché mantiene lo sforzo totale entro il limite e include più elementi ad alto rapporto valore/sforzo per me. È una selezione bilanciata che lascia ancora margine minimo di sforzo, quindi è difficilmente migliorabile senza rinunciare a oggetti molto efficienti. Se hai forti preferenze su uno degli elementi, sono disposto a valutare uno scambio puntuale purché restiamo entro il limite.'\}} \\
            }
        }
    }
    & & \\ \\

    \theutterance \stepcounter{utterance}  
    & & \multicolumn{4}{p{0.6\linewidth}}{
        \cellcolor[rgb]{0.9,0.9,0.9}{
            \makecell[{{p{\linewidth}}}]{
                \texttt{\tiny{[GM$\rangle$P2]}}
                \texttt{Stai partecipando a un gioco di negoziazione collaborativa.} \\
\\ 
\texttt{Insieme a un altro partecipante, dovete accordarvi su un unico set di oggetti da conservare. Ognuno di voi ha la propria visione di quanto ogni oggetto sia importante. Non sai come l'altro partecipante valuta gli oggetti. Inoltre, ti viene fornito lo sforzo che ogni oggetto richiede.} \\
\texttt{Potete accordarvi su un set solo se lo sforzo totale degli oggetti selezionati non supera un limite condiviso:} \\
\\ 
\texttt{LIMITE = 2984} \\
\\ 
\texttt{Ecco i valori di sforzo dei singoli oggetti:} \\
\\ 
\texttt{Sforzo oggetto = \{"A21": 350, "B93": 276, "A11": 616, "B33": 737, "C08": 531, "B39": 389, "B96": 24, "C85": 125, "A82": 338, "B31": 356, "B20": 143, "A60": 117, "B09": 257, "C81": 921, "B21": 789\}} \\
\\ 
\texttt{Ecco la tua visione personale sull'importanza di ogni oggetto:} \\
\\ 
\texttt{Valori di importanza dell'oggetto = \{"A11": 65, "B33": 121, "B21": 138, "C08": 262, "B31": 389, "B09": 461, "B93": 484, "A82": 508, "A60": 583, "B20": 668, "C81": 780, "B96": 783, "A21": 808, "C85": 822, "B39": 868\}} \\
\\ 
\texttt{Obiettivo:} \\
\\ 
\texttt{Il tuo obiettivo è negoziare un set condiviso di oggetti che ti avvantaggi il più possibile (cioè, massimizzi l'importanza totale per TE), rimanendo entro il LIMITEE. Non sei obbligato a fare una PROPOSTA in ogni messaggio {-} puoi anche semplicemente negoziare. Tutte le tattiche sono permesse!} \\
\\ 
\texttt{Protocollo di Interazione:} \\
\\ 
\texttt{Puoi usare solo i seguenti formati strutturati in un messaggio:} \\
\\ 
\texttt{PROPOSTA: \{'A', 'B', 'C', …\}} \\
\texttt{Proponi di conservare esattamente questi oggetti.} \\
\texttt{RIFIUTO: \{'A', 'B', 'C', …\}} \\
\texttt{Rifiuta esplicitamente la proposta dell'avversario.} \\
\texttt{ARGOMENTAZIONE: \{'...'\}} \\
\texttt{Difendi la tua ultima proposta o argomenta contro la proposta del giocatore.} \\
\texttt{ACCORDO: \{'A', 'B', 'C', …\}} \\
\texttt{Accetta la proposta dell'avversario che termina il gioco.} \\
\\ 
\\ 
\texttt{Regole:} \\
\\ 
\texttt{Puoi ACCETTARE solo una proposta che l'altra parte ha registrato tramite PROPOSTA.} \\
\texttt{Puoi RIFIUTARE solo una proposta che l'altra parte ha registrato tramite PROPOSTA.} \\
\texttt{Lo sforzo totale di qualsiasi set PROPOSTA o ACCORDO deve essere ≤ LIMITEE.} \\
\texttt{NON rivelare i tuoi punteggi di importanza nascosti.} \\
\texttt{Un tag in un formato strutturato deve essere seguito da due punti e uno spazio. L'argomento deve essere un set python contenente 0 o più stringhe.} \\
\texttt{Quindi, deve essere della forma TAG: \{...\}} \\
\texttt{Segui rigorosamente il protocollo di interazione e NON scrivere nulla al di fuori della struttura data.} \\
\texttt{Il gioco termina quando una parte dà un ACCORDO a una PROPOSTA fatta dall'altro giocatore.} \\
\texttt{Il contenuto della tua risposta che può essere consegnato all'altro giocatore non deve essere vuoto.} \\
\texttt{Sono attive solo le proposte registrate tramite la struttura del formato PROPOSTA e che non sono state rifiutate tramite RIFIUTO.} \\
\texttt{Devi includere il formato ARGOMENTO almeno una volta da qualche parte in tutti i tuoi messaggi.} \\
\\ 
\\ 
\texttt{Ora riceverai il primo messaggio dell'altro giocatore.} \\
\\ 
\texttt{PROPOSTA: \{'A11', 'B33', 'B31', 'B93', 'B09', 'A60'\}} \\
\\ 
\texttt{ARGOMENTAZIONE: \{'Propongo questo set perché mantiene lo sforzo totale entro il limite e include più elementi ad alto rapporto valore/sforzo per me. È una selezione bilanciata che lascia ancora margine minimo di sforzo, quindi è difficilmente migliorabile senza rinunciare a oggetti molto efficienti. Se hai forti preferenze su uno degli elementi, sono disposto a valutare uno scambio puntuale purché restiamo entro il limite.'\}} \\
            }
        }
    }
    & \\ \\

    \theutterance \stepcounter{utterance}  
    & & & \multicolumn{4}{p{0.6\linewidth}}{
        \cellcolor[rgb]{0.9,0.9,0.9}{
            \makecell[{{p{\linewidth}}}]{
                \texttt{\tiny{[GM$\langle$P2]}}
                \texttt{RIFIUTO: \{'A11', 'B33', 'B31', 'B93', 'B09', 'A60'\}} \\
\\ 
\texttt{ARGOMENTAZIONE: \{'La tua selezione è vicina al limite ma lascia fuori alcuni oggetti per me cruciali. Possiamo ottenere un set più efficiente per entrambi sostituendo elementi meno impattanti con altri a migliore rapporto valore/sforzo per me, restando entro il limite.'\}} \\
\\ 
\texttt{PROPOSTA: \{'A21', 'B39', 'C85', 'B96', 'A60', 'A82', 'B20'\}} \\
            }
        }
    }
     \\ \\

    \theutterance \stepcounter{utterance}  
    & & \multicolumn{4}{p{0.6\linewidth}}{
        \cellcolor[rgb]{0.9,0.9,0.9}{
            \makecell[{{p{\linewidth}}}]{
                \texttt{\tiny{[P1$\langle$GM]}}
                \texttt{RIFIUTO: \{'A11', 'B33', 'B31', 'B93', 'B09', 'A60'\}} \\
\\ 
\texttt{ARGOMENTAZIONE: \{'La tua selezione è vicina al limite ma lascia fuori alcuni oggetti per me cruciali. Possiamo ottenere un set più efficiente per entrambi sostituendo elementi meno impattanti con altri a migliore rapporto valore/sforzo per me, restando entro il limite.'\}} \\
\\ 
\texttt{PROPOSTA: \{'A21', 'B39', 'C85', 'B96', 'A60', 'A82', 'B20'\}} \\
            }
        }
    }
    & \\ \\

    \theutterance \stepcounter{utterance}  
    & \multicolumn{4}{p{0.6\linewidth}}{
        \cellcolor[rgb]{0.9,0.9,0.9}{
            \makecell[{{p{\linewidth}}}]{
                \texttt{\tiny{[P1$\rangle$GM]}}
                \texttt{RIFIUTO: \{'A21', 'B39', 'C85', 'B96', 'A60', 'A82', 'B20'\}} \\
\\ 
\texttt{ARGOMENTAZIONE: \{'La tua proposta supera il limite di sforzo: 350+389+125+24+117+338+143 = 1.486, che è sotto il limite, ma sacrifica troppa importanza relativa su oggetti chiave per me e lascia molto margine inutilizzato, rendendo il set sub{-}ottimale. Posso accogliere parte delle tue preferenze (B96, A60, B20, A82) ma vorrei includere alcuni elementi ad altissima efficienza per me. Propongo un compromesso che usa bene il budget e resta flessibile per un paio di scambi mirati.'\}} \\
\\ 
\texttt{PROPOSTA: \{'A11', 'B33', 'B31', 'B93', 'B09', 'A60', 'B96'\}} \\
            }
        }
    }
    & & \\ \\

    \theutterance \stepcounter{utterance}  
    & & \multicolumn{4}{p{0.6\linewidth}}{
        \cellcolor[rgb]{0.9,0.9,0.9}{
            \makecell[{{p{\linewidth}}}]{
                \texttt{\tiny{[GM$\rangle$P2]}}
                \texttt{RIFIUTO: \{'A21', 'B39', 'C85', 'B96', 'A60', 'A82', 'B20'\}} \\
\\ 
\texttt{ARGOMENTAZIONE: \{'La tua proposta supera il limite di sforzo: 350+389+125+24+117+338+143 = 1.486, che è sotto il limite, ma sacrifica troppa importanza relativa su oggetti chiave per me e lascia molto margine inutilizzato, rendendo il set sub{-}ottimale. Posso accogliere parte delle tue preferenze (B96, A60, B20, A82) ma vorrei includere alcuni elementi ad altissima efficienza per me. Propongo un compromesso che usa bene il budget e resta flessibile per un paio di scambi mirati.'\}} \\
\\ 
\texttt{PROPOSTA: \{'A11', 'B33', 'B31', 'B93', 'B09', 'A60', 'B96'\}} \\
            }
        }
    }
    & \\ \\

    \theutterance \stepcounter{utterance}  
    & & & \multicolumn{4}{p{0.6\linewidth}}{
        \cellcolor[rgb]{0.9,0.9,0.9}{
            \makecell[{{p{\linewidth}}}]{
                \texttt{\tiny{[GM$\langle$P2]}}
                \texttt{RIFIUTO: \{'A11', 'B33', 'B31', 'B93', 'B09', 'A60', 'B96'\}} \\
\\ 
\texttt{ARGOMENTAZIONE: \{'Capisco l’inclusione di B96 e A60. Tuttavia, il tuo set usa 2769 di sforzo (616+737+356+276+257+117+24) e sacrifica oggetti per me molto più impattanti a parità di sforzo. Possiamo aumentare l’efficienza complessiva mantenendo B96 e A60, aggiungendo alcuni miei pick chiave e restando ben entro il limite, così da lasciarci anche margine per un eventuale scambio puntuale.'\}} \\
\\ 
\texttt{PROPOSTA: \{'A21', 'B39', 'C85', 'B96', 'A60', 'A82', 'B20', 'B93'\}} \\
            }
        }
    }
     \\ \\

    \theutterance \stepcounter{utterance}  
    & & \multicolumn{4}{p{0.6\linewidth}}{
        \cellcolor[rgb]{0.9,0.9,0.9}{
            \makecell[{{p{\linewidth}}}]{
                \texttt{\tiny{[P1$\langle$GM]}}
                \texttt{RIFIUTO: \{'A11', 'B33', 'B31', 'B93', 'B09', 'A60', 'B96'\}} \\
\\ 
\texttt{ARGOMENTAZIONE: \{'Capisco l’inclusione di B96 e A60. Tuttavia, il tuo set usa 2769 di sforzo (616+737+356+276+257+117+24) e sacrifica oggetti per me molto più impattanti a parità di sforzo. Possiamo aumentare l’efficienza complessiva mantenendo B96 e A60, aggiungendo alcuni miei pick chiave e restando ben entro il limite, così da lasciarci anche margine per un eventuale scambio puntuale.'\}} \\
\\ 
\texttt{PROPOSTA: \{'A21', 'B39', 'C85', 'B96', 'A60', 'A82', 'B20', 'B93'\}} \\
            }
        }
    }
    & \\ \\

    \theutterance \stepcounter{utterance}  
    & \multicolumn{4}{p{0.6\linewidth}}{
        \cellcolor[rgb]{0.9,0.9,0.9}{
            \makecell[{{p{\linewidth}}}]{
                \texttt{\tiny{[P1$\rangle$GM]}}
                \texttt{RIFIUTO: \{'A21', 'B39', 'C85', 'B96', 'A60', 'A82', 'B20', 'B93'\}} \\
\\ 
\texttt{ARGOMENTAZIONE: \{'Apprezzo che tu mantenga B96 e A60 e aggiunga B93. Tuttavia, questa selezione lascia fuori troppi elementi con altissima resa per me (A11, B33, B31, B09) e impegna sforzo su oggetti a bassa resa per me (B39, C85). Propongo di mantenere B96, A60, B20 e inserire alcuni dei miei pick chiave. Sono disposto a rinunciare a uno tra B31 e B09 se ti serve spazio per un tuo elemento ad alta priorità, purché restiamo entro il limite.'\}} \\
\\ 
\texttt{PROPOSTA: \{'A11', 'B33', 'B31', 'B09', 'B93', 'A60', 'B96'\}} \\
            }
        }
    }
    & & \\ \\

    \theutterance \stepcounter{utterance}  
    & & \multicolumn{4}{p{0.6\linewidth}}{
        \cellcolor[rgb]{0.9,0.9,0.9}{
            \makecell[{{p{\linewidth}}}]{
                \texttt{\tiny{[GM$\rangle$P2]}}
                \texttt{RIFIUTO: \{'A21', 'B39', 'C85', 'B96', 'A60', 'A82', 'B20', 'B93'\}} \\
\\ 
\texttt{ARGOMENTAZIONE: \{'Apprezzo che tu mantenga B96 e A60 e aggiunga B93. Tuttavia, questa selezione lascia fuori troppi elementi con altissima resa per me (A11, B33, B31, B09) e impegna sforzo su oggetti a bassa resa per me (B39, C85). Propongo di mantenere B96, A60, B20 e inserire alcuni dei miei pick chiave. Sono disposto a rinunciare a uno tra B31 e B09 se ti serve spazio per un tuo elemento ad alta priorità, purché restiamo entro il limite.'\}} \\
\\ 
\texttt{PROPOSTA: \{'A11', 'B33', 'B31', 'B09', 'B93', 'A60', 'B96'\}} \\
            }
        }
    }
    & \\ \\

    \theutterance \stepcounter{utterance}  
    & & & \multicolumn{4}{p{0.6\linewidth}}{
        \cellcolor[rgb]{0.9,0.9,0.9}{
            \makecell[{{p{\linewidth}}}]{
                \texttt{\tiny{[GM$\langle$P2]}}
                \texttt{RIFIUTO: \{'A11', 'B33', 'B31', 'B09', 'B93', 'A60', 'B96'\}} \\
\\ 
\texttt{ARGOMENTAZIONE: \{'Riconosco il valore per te di A11, B33, B31, B09, ma questo set (616+737+356+257+276+117+24=2.383) consuma gran parte del budget lasciando fuori più elementi per me estremamente impattanti a parità di sforzo. Posso accogliere due tra i tuoi “alto rendimento” se contestualmente includiamo alcuni miei pick chiave più efficienti. Ti propongo un pacchetto che preserva B96, A60 e B93 come ponte tra le nostre preferenze, aggiunge due miei asset ad altissima resa e lascia ancora margine per un micro{-}swap mirato.'\}} \\
\\ 
\texttt{PROPOSTA: \{'A21', 'C85', 'B39', 'B96', 'A60', 'A82', 'B20', 'B93', 'B09'\}} \\
            }
        }
    }
     \\ \\

    \theutterance \stepcounter{utterance}  
    & & \multicolumn{4}{p{0.6\linewidth}}{
        \cellcolor[rgb]{0.9,0.9,0.9}{
            \makecell[{{p{\linewidth}}}]{
                \texttt{\tiny{[P1$\langle$GM]}}
                \texttt{RIFIUTO: \{'A11', 'B33', 'B31', 'B09', 'B93', 'A60', 'B96'\}} \\
\\ 
\texttt{ARGOMENTAZIONE: \{'Riconosco il valore per te di A11, B33, B31, B09, ma questo set (616+737+356+257+276+117+24=2.383) consuma gran parte del budget lasciando fuori più elementi per me estremamente impattanti a parità di sforzo. Posso accogliere due tra i tuoi “alto rendimento” se contestualmente includiamo alcuni miei pick chiave più efficienti. Ti propongo un pacchetto che preserva B96, A60 e B93 come ponte tra le nostre preferenze, aggiunge due miei asset ad altissima resa e lascia ancora margine per un micro{-}swap mirato.'\}} \\
\\ 
\texttt{PROPOSTA: \{'A21', 'C85', 'B39', 'B96', 'A60', 'A82', 'B20', 'B93', 'B09'\}} \\
            }
        }
    }
    & \\ \\

    \theutterance \stepcounter{utterance}  
    & \multicolumn{4}{p{0.6\linewidth}}{
        \cellcolor[rgb]{0.9,0.9,0.9}{
            \makecell[{{p{\linewidth}}}]{
                \texttt{\tiny{[P1$\rangle$GM]}}
                \texttt{RIFIUTO: \{'A21', 'C85', 'B39', 'B96', 'A60', 'A82', 'B20', 'B93', 'B09'\}} \\
\\ 
\texttt{ARGOMENTAZIONE: \{'Capisco il ponte su B96, A60, B93 e l’inclusione di B09. Tuttavia, stai allocando molto sforzo su elementi per me poco redditizi (A21, B39, C85) mentre lasci fuori asset con resa decisamente superiore per me. Ti propongo un compromesso che mantiene i tre punti di contatto (B96, A60, B93), include due miei pilastri (A11, B33) e ti lascia ancora spazio per un ulteriore innesto se necessario sostituendo un solo pezzo a bassa resa per me.'\}} \\
\\ 
\texttt{PROPOSTA: \{'A11', 'B33', 'B93', 'B09', 'A60', 'B96'\}} \\
            }
        }
    }
    & & \\ \\

    \theutterance \stepcounter{utterance}  
    & & \multicolumn{4}{p{0.6\linewidth}}{
        \cellcolor[rgb]{0.9,0.9,0.9}{
            \makecell[{{p{\linewidth}}}]{
                \texttt{\tiny{[GM$\rangle$P2]}}
                \texttt{RIFIUTO: \{'A21', 'C85', 'B39', 'B96', 'A60', 'A82', 'B20', 'B93', 'B09'\}} \\
\\ 
\texttt{ARGOMENTAZIONE: \{'Capisco il ponte su B96, A60, B93 e l’inclusione di B09. Tuttavia, stai allocando molto sforzo su elementi per me poco redditizi (A21, B39, C85) mentre lasci fuori asset con resa decisamente superiore per me. Ti propongo un compromesso che mantiene i tre punti di contatto (B96, A60, B93), include due miei pilastri (A11, B33) e ti lascia ancora spazio per un ulteriore innesto se necessario sostituendo un solo pezzo a bassa resa per me.'\}} \\
\\ 
\texttt{PROPOSTA: \{'A11', 'B33', 'B93', 'B09', 'A60', 'B96'\}} \\
            }
        }
    }
    & \\ \\

    \theutterance \stepcounter{utterance}  
    & & & \multicolumn{4}{p{0.6\linewidth}}{
        \cellcolor[rgb]{0.9,0.9,0.9}{
            \makecell[{{p{\linewidth}}}]{
                \texttt{\tiny{[GM$\langle$P2]}}
                \texttt{RIFIUTO: \{'A11', 'B33', 'B93', 'B09', 'A60', 'B96'\}} \\
\\ 
\texttt{ARGOMENTAZIONE: \{'Posso convergere su B96, A60 e B93, ma tenere insieme A11, B33 e B09 blocca troppo budget su elementi che per me hanno resa inferiore rispetto ad alternative. Ti propongo un compromesso che mantiene i tre punti di contatto, aggiunge due miei asset ad altissima efficienza e lascia ancora margine: se vuoi, possiamo scambiare uno tra A21/C85 con uno dei tuoi pilastri in un passo successivo pur restando entro il limite.'\}} \\
\\ 
\texttt{PROPOSTA: \{'B96', 'A60', 'B93', 'A21', 'C85', 'B20', 'A82'\}} \\
            }
        }
    }
     \\ \\

    \theutterance \stepcounter{utterance}  
    & & \multicolumn{4}{p{0.6\linewidth}}{
        \cellcolor[rgb]{0.9,0.9,0.9}{
            \makecell[{{p{\linewidth}}}]{
                \texttt{\tiny{[P1$\langle$GM]}}
                \texttt{RIFIUTO: \{'A11', 'B33', 'B93', 'B09', 'A60', 'B96'\}} \\
\\ 
\texttt{ARGOMENTAZIONE: \{'Posso convergere su B96, A60 e B93, ma tenere insieme A11, B33 e B09 blocca troppo budget su elementi che per me hanno resa inferiore rispetto ad alternative. Ti propongo un compromesso che mantiene i tre punti di contatto, aggiunge due miei asset ad altissima efficienza e lascia ancora margine: se vuoi, possiamo scambiare uno tra A21/C85 con uno dei tuoi pilastri in un passo successivo pur restando entro il limite.'\}} \\
\\ 
\texttt{PROPOSTA: \{'B96', 'A60', 'B93', 'A21', 'C85', 'B20', 'A82'\}} \\
            }
        }
    }
    & \\ \\

    \theutterance \stepcounter{utterance}  
    & \multicolumn{4}{p{0.6\linewidth}}{
        \cellcolor[rgb]{0.9,0.9,0.9}{
            \makecell[{{p{\linewidth}}}]{
                \texttt{\tiny{[P1$\rangle$GM]}}
                \texttt{RIFIUTO: \{'B96', 'A60', 'B93', 'A21', 'C85', 'B20', 'A82'\}} \\
\\ 
\texttt{ARGOMENTAZIONE: \{'Accolgo i tre punti di contatto (B96, A60, B93) ma A21 e C85 per me hanno resa bassa rispetto ad alternative. Possiamo tenere il tuo asse efficiente B20 e A82 se inseriamo almeno un mio pilastro. Propongo uno scambio mirato: sostituiamo A21 e C85 con A11 e B33. Restiamo entro il limite, alziamo molto la resa per me e manteniamo diversi tuoi pick chiave.'\}} \\
\\ 
\texttt{PROPOSTA: \{'A11', 'B33', 'B93', 'B20', 'A82', 'A60', 'B96'\}} \\
            }
        }
    }
    & & \\ \\

    \theutterance \stepcounter{utterance}  
    & & \multicolumn{4}{p{0.6\linewidth}}{
        \cellcolor[rgb]{0.9,0.9,0.9}{
            \makecell[{{p{\linewidth}}}]{
                \texttt{\tiny{[GM$\rangle$P2]}}
                \texttt{RIFIUTO: \{'B96', 'A60', 'B93', 'A21', 'C85', 'B20', 'A82'\}} \\
\\ 
\texttt{ARGOMENTAZIONE: \{'Accolgo i tre punti di contatto (B96, A60, B93) ma A21 e C85 per me hanno resa bassa rispetto ad alternative. Possiamo tenere il tuo asse efficiente B20 e A82 se inseriamo almeno un mio pilastro. Propongo uno scambio mirato: sostituiamo A21 e C85 con A11 e B33. Restiamo entro il limite, alziamo molto la resa per me e manteniamo diversi tuoi pick chiave.'\}} \\
\\ 
\texttt{PROPOSTA: \{'A11', 'B33', 'B93', 'B20', 'A82', 'A60', 'B96'\}} \\
            }
        }
    }
    & \\ \\

    \theutterance \stepcounter{utterance}  
    & & & \multicolumn{4}{p{0.6\linewidth}}{
        \cellcolor[rgb]{0.9,0.9,0.9}{
            \makecell[{{p{\linewidth}}}]{
                \texttt{\tiny{[GM$\langle$P2]}}
                \texttt{RIFIUTO: \{'A11', 'B33', 'B93', 'B20', 'A82', 'A60', 'B96'\}} \\
\\ 
\texttt{ARGOMENTAZIONE: \{'Capisco la tua esigenza di inserire A11 e B33, ma questo pacchetto blocca troppo del budget su elementi che per me rendono meno di alternative molto efficienti per me. Possiamo però costruire un compromesso equilibrato mantenendo i tre punti di contatto (B96, A60, B93), conservando il tuo asse B20, e scambiando uno solo tra i tuoi pilastri con due miei pick chiave ad altissima resa. Ti lascio margine per un eventuale micro{-}swap successivo pur restando ampiamente sotto il limite.'\}} \\
\\ 
\texttt{PROPOSTA: \{'B96', 'A60', 'B93', 'B20', 'A82', 'A21', 'B39'\}} \\
            }
        }
    }
     \\ \\

    \theutterance \stepcounter{utterance}  
    & & \multicolumn{4}{p{0.6\linewidth}}{
        \cellcolor[rgb]{0.9,0.9,0.9}{
            \makecell[{{p{\linewidth}}}]{
                \texttt{\tiny{[P1$\langle$GM]}}
                \texttt{RIFIUTO: \{'A11', 'B33', 'B93', 'B20', 'A82', 'A60', 'B96'\}} \\
\\ 
\texttt{ARGOMENTAZIONE: \{'Capisco la tua esigenza di inserire A11 e B33, ma questo pacchetto blocca troppo del budget su elementi che per me rendono meno di alternative molto efficienti per me. Possiamo però costruire un compromesso equilibrato mantenendo i tre punti di contatto (B96, A60, B93), conservando il tuo asse B20, e scambiando uno solo tra i tuoi pilastri con due miei pick chiave ad altissima resa. Ti lascio margine per un eventuale micro{-}swap successivo pur restando ampiamente sotto il limite.'\}} \\
\\ 
\texttt{PROPOSTA: \{'B96', 'A60', 'B93', 'B20', 'A82', 'A21', 'B39'\}} \\
            }
        }
    }
    & \\ \\

    \theutterance \stepcounter{utterance}  
    & \multicolumn{4}{p{0.6\linewidth}}{
        \cellcolor[rgb]{0.9,0.9,0.9}{
            \makecell[{{p{\linewidth}}}]{
                \texttt{\tiny{[P1$\rangle$GM]}}
                \texttt{RIFIUTO: \{'B96', 'A60', 'B93', 'B20', 'A82', 'A21', 'B39'\}} \\
\\ 
\texttt{ARGOMENTAZIONE: \{'Apprezzo che tu mantenga B96, A60, B93 e B20. Tuttavia A21 e soprattutto B39 per me hanno resa bassa rispetto ad alternative, mentre ho bisogno di includere almeno un mio pilastro ad altissima resa. Per avvicinarci, ti propongo di sostituire A21 e B39 con A11 e B33: così alziamo molto l’efficienza per me mantenendo quattro dei tuoi elementi preferiti. Se ti serve un aggiustamento successivo, posso valutare di scambiare B09 con uno dei miei pilastri in un secondo passo.'\}} \\
\\ 
\texttt{PROPOSTA: \{'A11', 'B33', 'B93', 'B20', 'A82', 'A60', 'B96'\}} \\
            }
        }
    }
    & & \\ \\

    \theutterance \stepcounter{utterance}  
    & & \multicolumn{4}{p{0.6\linewidth}}{
        \cellcolor[rgb]{0.9,0.9,0.9}{
            \makecell[{{p{\linewidth}}}]{
                \texttt{\tiny{[GM$\rangle$P2]}}
                \texttt{RIFIUTO: \{'B96', 'A60', 'B93', 'B20', 'A82', 'A21', 'B39'\}} \\
\\ 
\texttt{ARGOMENTAZIONE: \{'Apprezzo che tu mantenga B96, A60, B93 e B20. Tuttavia A21 e soprattutto B39 per me hanno resa bassa rispetto ad alternative, mentre ho bisogno di includere almeno un mio pilastro ad altissima resa. Per avvicinarci, ti propongo di sostituire A21 e B39 con A11 e B33: così alziamo molto l’efficienza per me mantenendo quattro dei tuoi elementi preferiti. Se ti serve un aggiustamento successivo, posso valutare di scambiare B09 con uno dei miei pilastri in un secondo passo.'\}} \\
\\ 
\texttt{PROPOSTA: \{'A11', 'B33', 'B93', 'B20', 'A82', 'A60', 'B96'\}} \\
            }
        }
    }
    & \\ \\

    \theutterance \stepcounter{utterance}  
    & & & \multicolumn{4}{p{0.6\linewidth}}{
        \cellcolor[rgb]{0.9,0.9,0.9}{
            \makecell[{{p{\linewidth}}}]{
                \texttt{\tiny{[GM$\langle$P2]}}
                \texttt{RIFIUTO: \{'A11', 'B33', 'B93', 'B20', 'A82', 'A60', 'B96'\}} \\
\\ 
\texttt{ARGOMENTAZIONE: \{'Vedo il tuo bisogno di includere almeno un tuo pilastro, ma tenere insieme sia A11 che B33 vincola troppo budget su elementi per me meno efficienti di alternative. Sono disposto a includere UNO dei tuoi pilastri ad altissima resa se contestualmente manteniamo i nostri tre punti di contatto (B96, A60, B93) e il mio asse efficiente (B20, A82). Questo crea un set equilibrato, ampiamente sotto il limite, e lascia ancora margine per un micro{-}swap successivo se necessario.'\}} \\
\\ 
\texttt{PROPOSTA: \{'A11', 'B33', 'B93', 'B20', 'A82', 'A60', 'B96'\}} \\
            }
        }
    }
     \\ \\

    \theutterance \stepcounter{utterance}  
    & & \multicolumn{4}{p{0.6\linewidth}}{
        \cellcolor[rgb]{0.9,0.9,0.9}{
            \makecell[{{p{\linewidth}}}]{
                \texttt{\tiny{[P1$\langle$GM]}}
                \texttt{RIFIUTO: \{'A11', 'B33', 'B93', 'B20', 'A82', 'A60', 'B96'\}} \\
\\ 
\texttt{ARGOMENTAZIONE: \{'Vedo il tuo bisogno di includere almeno un tuo pilastro, ma tenere insieme sia A11 che B33 vincola troppo budget su elementi per me meno efficienti di alternative. Sono disposto a includere UNO dei tuoi pilastri ad altissima resa se contestualmente manteniamo i nostri tre punti di contatto (B96, A60, B93) e il mio asse efficiente (B20, A82). Questo crea un set equilibrato, ampiamente sotto il limite, e lascia ancora margine per un micro{-}swap successivo se necessario.'\}} \\
\\ 
\texttt{PROPOSTA: \{'A11', 'B33', 'B93', 'B20', 'A82', 'A60', 'B96'\}} \\
            }
        }
    }
    & \\ \\

    \theutterance \stepcounter{utterance}  
    & \multicolumn{4}{p{0.6\linewidth}}{
        \cellcolor[rgb]{0.9,0.9,0.9}{
            \makecell[{{p{\linewidth}}}]{
                \texttt{\tiny{[P1$\rangle$GM]}}
                \texttt{RIFIUTO: \{'A11', 'B33', 'B93', 'B20', 'A82', 'A60', 'B96'\}} \\
\\ 
\texttt{ARGOMENTAZIONE: \{'Accolgo la tua apertura a includere uno dei miei pilastri. Per convergere, manteniamo i tre punti di contatto (B96, A60, B93) e il tuo asse efficiente (B20, A82), inserendo UN solo mio pilastro ad altissima resa. Questo resta ben sotto il limite e lascia spazio per un eventuale micro{-}swap.'\}} \\
\\ 
\texttt{PROPOSTA: \{'A11', 'B93', 'B20', 'A82', 'A60', 'B96'\}} \\
            }
        }
    }
    & & \\ \\

    \theutterance \stepcounter{utterance}  
    & & \multicolumn{4}{p{0.6\linewidth}}{
        \cellcolor[rgb]{0.9,0.9,0.9}{
            \makecell[{{p{\linewidth}}}]{
                \texttt{\tiny{[GM$\rangle$P2]}}
                \texttt{RIFIUTO: \{'A11', 'B33', 'B93', 'B20', 'A82', 'A60', 'B96'\}} \\
\\ 
\texttt{ARGOMENTAZIONE: \{'Accolgo la tua apertura a includere uno dei miei pilastri. Per convergere, manteniamo i tre punti di contatto (B96, A60, B93) e il tuo asse efficiente (B20, A82), inserendo UN solo mio pilastro ad altissima resa. Questo resta ben sotto il limite e lascia spazio per un eventuale micro{-}swap.'\}} \\
\\ 
\texttt{PROPOSTA: \{'A11', 'B93', 'B20', 'A82', 'A60', 'B96'\}} \\
            }
        }
    }
    & \\ \\

    \theutterance \stepcounter{utterance}  
    & & & \multicolumn{4}{p{0.6\linewidth}}{
        \cellcolor[rgb]{0.9,0.9,0.9}{
            \makecell[{{p{\linewidth}}}]{
                \texttt{\tiny{[GM$\langle$P2]}}
                \texttt{ARGOMENTAZIONE: \{'Apprezzo la tua flessibilità e la struttura proposta: mantiene i tre punti di contatto e il tuo asse efficiente, includendo uno dei tuoi pilastri. È un compromesso coerente con quanto discusso e lascia margine utile.'\}} \\
\\ 
\texttt{ACCORDO: \{'A11', 'B93', 'B20', 'A82', 'A60', 'B96'\}} \\
            }
        }
    }
     \\ \\

    \theutterance \stepcounter{utterance}  
    & & & \multicolumn{2}{p{0.3\linewidth}}{
        \cellcolor[rgb]{0.9,0.9,0.9}{
            \makecell[{{p{\linewidth}}}]{
                \texttt{\tiny{[GM$|$GM]}}
                \texttt{game successful} \\
            }
        }
    }
    & & \\ \\

    \theutterance \stepcounter{utterance}  
    & & & \multicolumn{2}{p{0.3\linewidth}}{
        \cellcolor[rgb]{0.9,0.9,0.9}{
            \makecell[{{p{\linewidth}}}]{
                \texttt{\tiny{[GM$|$GM]}}
                \texttt{end game} \\
            }
        }
    }
    & & \\ \\

\end{supertabular}
}

\end{document}
